\section{Returning the polarized representation to the original CM field}

% BEGIN D70W1
\paragraph{D70W1. Descend a cusp through one cyclic step with the prescribed Galois representation.}
Let $E/F$ be a cyclic extension of prime degree between CM fields. Fix
$\iota:\overline{\mathbb Q}_\ell\simeq\mathbb C$ and a continuous
$n$-dimensional representation $r$ of $G_F$ whose restriction to
$G_E$ is absolutely irreducible. Suppose a regular algebraic cusp
$\Pi$ of $\mathrm{GL}_n(\mathbb A_E)$ is attached to that restriction
in the sense of $\mathrm{AttC}$. Thus at almost every $w$ the
geometric Frobenius roots are the unitary-normalized Satake multiset
multiplied by $q_w^{(n-1)/2}$, with coefficients compared by $\iota$.

For each generator of the deck group the conjugate cusp has the
same Satake multisets at almost every place. Indeed, the corresponding
Frobenius elements are conjugate under $G_F$ and $r$ extends there.
Strong multiplicity one identifies the two cusps. Cyclic cuspidal
descent (Arthur--Clozel, \emph{Simple algebras, base change, and
the advanced theory of the trace formula}, Ch.~III) gives a cusp
$\pi_0$ over $F$ with base change $\Pi$. Every archimedean local
extension here is $\mathbb C/\mathbb C$. Its infinitesimal character
is unchanged, so the exponents at each embedding downstairs are
distinct elements of $(n-1)/2+\mathbb Z$. This is exactly regular
algebraicity in the source.

Use Harris--Lan--Taylor--Thorne, \emph{On the rigid cohomology of
certain Shimura varieties}, Thm.~A (Cor.~7.14),
doi:10.1186/s40687-016-0078-5, to attach the semisimple
$s_{\pi_0}$ to this regular algebraic cusp over a CM field. At good
places base change raises the Satake roots to the residue degree
$f$, and $q_w^{(n-1)/2}=(q_v^{(n-1)/2})^f$. Chebotarev density and
Brauer--Nesbitt therefore identify $(s_{\pi_0}|_{G_E})^{\rm ss}$
with $r|_{G_E}$. Since the latter is irreducible, the restriction
itself is isomorphic to it. Identify the two on $G_E$. For $g\in G_F$
the matrix $s_{\pi_0}(g)^{-1}r(g)$ now commutes with $r(G_E)$ by
normality. Schur's lemma makes it a scalar $\eta(g)$. Multiplying
the matrices shows that $\eta$ is a character trivial on $G_E$.
Via reciprocity take the finite order Hecke character matching
$\eta$ at geometric Frobenius in the attachment convention (CIP1).
Twisting $\pi_0$ by it produces a cusp attached to the given $r$.
The twist pulls back trivially to $E$ and has trivial infinitesimal
character. This proves the step before any polarization of the
descendant has been used. Rank one also follows directly from
global reciprocity.
% END D70W1

% BEGIN D70W2
\paragraph{D70W2. Solvable descent and the conjugate-dual identity.}
Take the fields $K\subset L$ and the data of $\mathrm{DownCsdKAlt}$.
Choose a subnormal series of the finite solvable group
$\mathrm{Gal}(L/K)$ with successive quotients cyclic of prime order.
All intermediate fields are CM. Indeed the normal closure of $L$ is
a compositum of CM conjugates, with central complex conjugation.
It stabilizes every intermediate field. Its fixed field there is
totally real, and its action is nontrivial since the intermediate
field contains $K$. On each rung $r$ is irreducible because it is
irreducible on $G_L$. Starting with the cusp in the upper
$\mathrm{AutCsd0}$, apply D70W1 down this tower. This also covers the
degree-one case by taking the original cusp. We obtain a regular
algebraic cusp $\pi$ attached to the original $r$ over $K$.

Write $c$ for the specified complex conjugation. On geometric
Frobenius at a good place $v$, the multiplier
$\varepsilon^{1-n}$ in the given pairing is $q_v^{n-1}$.
Thus the roots of $r$ at $c v$ are $q_v^{n-1}$ times the inverses
of the roots at $v$. Dividing out the two factors
$q_v^{(n-1)/2}$ in the comparison gives
$\mathrm{Sat}(\pi_{cv})=\mathrm{Sat}(\pi_v)^{-1}$ as multisets.
Strong multiplicity one identifies $\pi^c$ with $\pi^\vee$.
We have proved the entire lower $\mathrm{AutCsd0}$. Also, any
regular algebraic cusp attached to this same $r$ has the same Satake
multisets almost everywhere and hence the same local components.
This observation deals with the universal cusp quantifier in
$\mathrm{WitA}$.
% END D70W2

% BEGIN D70W3
\paragraph{D70W3. The support downstairs consists of characters.}
For a regular algebraic conjugate self-dual cusp over a CM field,
local--global compatibility identifies the semisimplified underlying
Weil representation of its algebraically normalized local parameter
with that of its attached Galois representation. Away from $\ell$
use Varma, \emph{Forum Math.\ Sigma} 12 (2024), e21,
doi:10.1017/fms.2024.7. At places over $\ell$ use
Barnet-Lamb--Gee--Geraghty--Taylor, \emph{Local--global
compatibility for $l=p$, II}, doi:10.24033/asens.2212, which gives
semisimplified compatibility for every regular weight. Here the
Galois representation is potentially semistable at $\ell$.
The corresponding results in rank one are local reciprocity.
The comparisons apply to the actual $r$ and $r|_{G_L}$ by their
cofinite attachments, irreducibility and Chebotarev. In the local
Langlands correspondence for $\mathrm{GL}_n$, the irreducible Weil
constituents, with multiplicities and the exponents in their
segments, give the supercuspidal support of normalized induction.
In particular, a sum of Weil characters gives toral support. This
uses the usual compatibility of local Langlands with the
Langlands classification. Passing from the algebraic normalization
to the local cusp normalization removes one common unramified twist.

At $v\nmid\ell$, $\mathrm{IneAb}$ supplies continuous characters
of the \emph{whole} decomposition group whose restrictions give the
characteristic polynomials of $r$ on inertia. Each such character
has finite image on inertia, by local reciprocity: local units
are finite modulo a pro-$p$ group, where $p\ne\ell$, and an
$\ell$-adic character on that pro-$p$ group has finite image.
One may take a finite coefficient field by C69C1.
The Weil--Deligne construction removes a unipotent exponential
from the inertia action. On inertia the exponential commutes with
the finite Weil action, hence leaves its eigenvalues unchanged.
The traces on the resulting common finite inertia quotient identify
the Weil inertia representation with the sum of the specified
inertial characters. Frobenius preserves each isotypic space
because each character extends to the decomposition group. There
inertia acts by scalars, and triangularizing Frobenius gives only
one-dimensional Weil constituents.

At $v\mid\ell$ take the full local flag supplied by $\mathrm{Flagged}$.
Its factors are characters of the bottom decomposition group.
Potentially semistable representations are closed under subquotients.
After taking a common finite extension of semistability the
semistable period functor is exact on this category, with descent
data. The resulting Weil--Deligne functor thus sends the flag to
a filtration with one-dimensional Weil factors on the bottom
Weil group. In both cases the local cusp has $\mathrm{TorAt}$.
No determination of the monodromy operator is needed for this
support calculation.
% END D70W3

% BEGIN D70W4
\paragraph{D70W4. Potential unlinking and the common center of the local bounds descend.}
At $v\mid\ell$, for characters $\kappa_i$ of a local flag let
$\delta_{ij}=\kappa_i/(\kappa_j\varepsilon)$. The source's
$\mathrm{PotUnl}$ requires that every $\delta_{ij}$ have infinite
order, including repeated indices, at every prolongation of $v$.
Restriction of a full flag to a local subgroup preserves its
one-dimensional factors, and Jordan--H\"older gives the same
multiset for any other full flag there. A character has finite
order exactly when its restriction to a subgroup of finite index
has finite order. Indeed a power trivial on the subgroup factors
through a finite quotient, by taking its normal core. Consequently
lower $\mathrm{PotUnl}$ implies upper $\mathrm{PotUnl}$ at every
$w\mid v$. Conversely, if the upper condition holds at all $w\mid v$,
take any lower prolongation and restrict its flag to the upper
decomposition group there. The upper condition forces each lower
ratio to have infinite order. If $\mathrm{Unl1}$ holds below,
$\mathrm{Keeps}$ supplies $\mathrm{Unl1}$ above at each such $w$,
and the first clause of the upper $\mathrm{WitA}$ yields
$\mathrm{PotUnl}$ there. This proves the first lower clause.

Apply the second upper clause to the chosen upper cusp $\Pi$ and
fix its real center $t$ once. Given a finite $v$ with
$v\nmid\ell$ or lower $\mathrm{PotUnl}$, choose $w\mid v$. The preceding
argument makes $w$ eligible. For any support character $\xi$ of
$\pi_v$, its pullback by the local norm occurs in the supercuspidal
support of $\Pi_w$. To see this directly, apply local--global
compatibility from D70W3 at both ends. Restriction of a
Weil--Deligne representation to a finite extension is compatible
with restriction of the Galois representation (at $\ell$ use the
same potentially semistable period module with restricted descent
data). Restrict the filtration giving its Weil character factors.
The common unramified normalization also restricts, since
$q_w=q_v^{f(w/v)}$. Local reciprocity now identifies the restricted
character with $\xi\circ N_{L_w/K_v}$.
A uniformizer above has norm a unit times $\varpi_v^f$, where
$f=f(w/v)$, and a smooth character has unit modulus on units.
The upper inequality is exactly
\[
 q_v^{f(t-1/2)}< |\xi(\varpi_v)|^f <q_v^{f(t+1/2)} .
\]
Taking positive $f$th roots keeps both inequalities strict.
Together with D70W3 this gives $\mathrm{WinAt}(\pi,t,v)$ for all
eligible $v$. D70W2 allows every regular algebraic attached cusp
in the lower clause. The same choice of $t$ works throughout.
% END D70W4

% BEGIN D70W5
\paragraph{D70W5. The full polarized solvable descent on the marked representation.}
For every $n\ge1$ and every finite solvable Galois extension of CM
fields in the original $\mathrm{DownCsdKAlt}$, D70W2 supplies the
lower $\mathrm{AutCsd0}$ and D70W3--4 its witness $\mathrm{WitA}$
under the stated continuity assumption. These are the two fields of
$\mathrm{IsAutCsdA}$. All comparisons used the input $\iota$, the
given $\ell$ and the given $r$. The finite characters introduced on
cyclic steps were chosen precisely to recover $r$.
In particular this result applies after the inverse-cyclotomic
rank-three twist on the marked Gaussian host of M65A4.
It supplies the descent input of L69D4 at that mark with its
full finite-place witness.
% END D70W5
